\section{Methods}
\label{sec:methods}

This section states each algorithm we evaluate, the choices we made, and -- in
\textbf{Why-it-was-chosen} blocks -- the mechanism we expect to see fire on
\mimic{}'s heterogeneity. The actual firing or non-firing is the entire job
of Section~\ref{sec:mechanism}.

\subsection{\fedavg{}}
The reference algorithm. Each round, $C \!=\! 5$ clients are sampled uniformly
without replacement, each runs $E \!=\! 1$ local epoch with batch size $256$ and
Adam, and the server averages the resulting weights with size weights. Code:
\texttt{flopt/fedavg.py}.

\textbf{Why-it-was-chosen.} \fedavg{} is the floor of any sensible federated
ablation: any modification has to beat it. The implementation deliberately uses
size weighting (not uniform), so it inherits the data imbalance directly --
which is what we want, because we want to see that imbalance get expressed as
a fairness failure that subsequent algorithms then repair.

\subsection{\fedprox{} (\texttt{flopt/fedprox.py})}
\fedprox{} adds a proximal term to each client's local objective:
\[
F_k^{\mathrm{prox}}(w; w_t) \;=\; F_k(w) \;+\; \frac{\mu}{2}\, \|w - w_t\|_2^2,
\]
where $w_t$ is the global model received at the start of round $t$. We sweep
$\mu \in \{0,\,10^{-3},\,10^{-2},\,10^{-1}\}$. Aggregation is unchanged from
\fedavg{}; the only modification is per-client local training. Code excerpt:
\begin{lstlisting}
prox_term = 0.0
for n, p in model.named_parameters():
    prox_term = prox_term + (mu / 2.0) * ((p - global_state[n]) ** 2).sum()
loss = ce(logits, y) + prox_term
\end{lstlisting}
Iteration is over \texttt{named\_parameters()} rather than
\texttt{state\_dict().items()} so buffers (running mean/var) do not appear in
the proximal term.

\textbf{Why-it-was-chosen.} On heterogeneous data, the per-client minima
$\{w_k^*\}$ are far apart; a $E\!=\!1$ local pass already pulls each $w_{t,k}$
toward $w_k^*$, far from the consensus, so the size-weighted average has high
variance. The proximal term is a quadratic anchor that costs the local loss
$O(\|w_{t,k} - w_t\|^2)$, which is exactly the quantity \fedavg{} pays for in
the next round in lost agreement. The cost-of-disagreement is paid \emph{once},
inside each client, instead of accumulated globally.

The expected mechanism on this dataset is: (i) lower per-round drift norms in
small Neuro units; (ii) faster convergence in number of rounds; (iii) higher
worst-client recall, because the small units no longer get over-corrected
toward MICU. Section~\ref{sec:mechanism-fedprox} verifies this on the actual
drift telemetry.

\subsection{\cvar{}-style aggregation reweighting}
The \cvar{} variant we implement does not modify local training; it modifies
\emph{aggregation}. After a round, sort selected clients by their local loss
$\hat{F}_{k}$ (highest loss = worst client). Replace the size-weighted
aggregator with a CVaR-$\alpha$ aggregator:
\[
w_{t+1} \;=\; \frac{1}{|S_t^{\alpha}|} \sum_{k \in S_t^{\alpha}} w_{t,k}^{(E)},
\qquad S_t^{\alpha} = \{k \in S_t : \hat{F}_k \geq \mathrm{Quantile}_\alpha(\hat{F})\}.
\]
$\alpha \!=\! 0$ recovers uniform-over-selected; $\alpha \!=\! 1$ keeps only the
single worst client. We sweep $\alpha \in \{0, 0.5, 0.75, 0.9, 0.95\}$.

\textbf{Why-it-was-chosen.} Risk-averse statistical learning argues that
\cvar{}-of-loss is the right object when you care about the worst customer
(here, ICU). Under perfect selection, increasing $\alpha$ should monotonically
improve worst-client recall at the price of accuracy.

\textbf{The thing we expected, and what actually happens.} With $C\!=\!5$
sampled clients per round, the tail set $S_t^{\alpha}$ is
$\lfloor (1-\alpha)\!\cdot\!C \rfloor$ clients. For $\alpha \!=\! 0.9$ this is
$\lfloor 0.5 \rfloor \!=\! 0$, and the implementation falls back to ``keep the
single worst.'' For $\alpha \!=\! 0.95$ this is also $0$, falling back the same
way. So \cvar{}-$0.9$, \cvar{}-$0.95$, and even \cvar{}-$0.75$ ($\lfloor 1.25
\rfloor = 1$) end up with $1$-or-$2$ client aggregation, which is identical in
expectation. Section~\ref{sec:mechanism-cvar} returns to this discreteness.

\subsection{Centralized and local-only baselines}
\textbf{Centralized}: pool all $54{,}853$ training rows, train the same
TabularMLP for $50$ epochs with the same Adam, evaluate on each client's test
fold. Provides the per-client \emph{ceiling}.

\textbf{Local-only}: train a fresh model on each client's training set,
evaluate only on that client's test fold. Aggregate by recording per-client
metrics (no global model). Provides the per-client \emph{floor}.

\textbf{Why-they-were-chosen.} Centralized is unattainable in deployment but
sets the upper envelope for what any federated method could possibly do.
Local-only is the literal alternative to federation if no data sharing happens.
Their gap is the room federation has to operate in.

\subsection{Hyperparameter search: grid \& differential evolution}
We search three knobs: $E \!\in\! \{1,2,3\}$ local epochs, $C \!\in\!
\{3,4,\dots,9\}$ clients per round, $\eta$ learning rate. For grid search we
use the discretisation $E \!\in\! \{1,2,3\}$, $C \!\in\! \{5,10,15\}$,
$\eta \!\in\! \{0.003,0.005,0.01\}$ -- $27$ combinations.

For DE search we use \texttt{scipy.optimize.differential\_evolution} with
\texttt{popsize}=$15$, \texttt{maxiter}=$15$, fitness $= \text{loss} + \beta_C
\!\cdot\! \text{(comm/}10^9)$, $\beta_C = 1$, evaluating each candidate on a
held-out evaluation subset of $30$ clients ($K\!=\!9$ for the natural-clients
case, $K \!=\! 30$ for Dirichlet). The proposal-alignment run is at $80$ rounds
per evaluation; the original run is at $30$. The proposal-alignment run uses
$90$ DE evaluations; the original uses $48$.

\textbf{Why-it-was-chosen.} Grid search at three levels per axis is the cheapest
sanity baseline. DE is the cheapest stochastic search that reliably handles
mixed integer/continuous spaces (we round $E$ and $C$ to ints inside the
evaluator). The hypothesis is that DE will find an off-grid $\eta$ that grid
misses; Section~\ref{sec:mechanism-search} confirms this exactly.

\subsection{LP duality for communication shadow prices}
We formulate a static policy program: choose a single configuration $(E, C,
\eta)$ subject to a total-comm budget. Let $\mathcal{X}$ be the candidate set
(the GA history, $48$/$90$ candidates, scored by loss and bytes). Solve
\[
\min_{p \in \Delta^{|\mathcal{X}|}} \;\sum_x p_x \!\cdot\! \mathrm{loss}_x
\quad \text{s.t.} \quad \sum_x p_x \!\cdot\! \mathrm{bytes}_x \leq B,
\]
which is a $|\mathcal{X}|$-dimensional simplex LP with one extra inequality. The
dual variable $\lambda$ on the budget constraint is the
\emph{communication shadow price}: the marginal loss reduction per extra byte
of budget. We solve at $12$ budgets log-spaced between $\sim 1.6\!\times\!10^{8}$
and $\sim 8.7\!\times\!10^{8}$ for the dense (MLP) program, and from
$\sim 2.6\!\times\!10^{6}$ for the sparse-aware (logistic) program.

We log primal feasibility, dual feasibility, complementary slackness, and a
stationarity residual; the full diagnostic is in
\texttt{outputs/full\_mimic\_iv\_proposal\_alignment/lp/kkt\_diagnostics.csv}.
All $24$ programs solve to optimal with KKT residuals below $10^{-9}$.

\textbf{Why-it-was-chosen.} Doing this LP for the headline run answers a
question that no individual method does: ``would another byte of communication
have helped?'' If $\lambda$ is huge, yes -- buy more. If $\lambda \!\approx\! 0$,
no -- spend the budget elsewhere (model size, search budget, calibration).
On \mimic{} $\lambda$ collapses 13 orders of magnitude as soon as the budget
exits the active region; that is the LP saying ``stop spending bytes.''

\subsection{Sparsity-aware communication}
Top-$k$ compression: at the end of each local epoch, keep only the largest
$k\!=\!\{1, 5, 10, 25, 100\}\%$ of update entries by absolute value, zero the rest,
upload only the survivors. We measure both the natural sparsity (fraction of
entries that would be zero \emph{without} truncation) and the achieved
compression ratio (sparse bytes / dense bytes). The summary is in
\texttt{outputs/full\_mimic\_iv\_proposal\_alignment/lp/sparsity\_summary.csv}.

\textbf{Why-it-was-chosen.} Two questions: (a) \emph{do} updates have natural
sparsity to exploit? (b) does forced top-$k$ trade lossy compression for
budget? Spoiler: (a) no, $95.8$--$97.7\%$ of entries are nonzero
\emph{without} truncation; (b) yes, top-$1\%$ achieves $\sim$$49\!\times$
compression. So the compression on \mimic{} is mechanistically lossy
truncation, not a free natural-sparsity dividend.

\subsection{Dirichlet study (synthetic non-IID stress test)}
The natural-clients case has $K \!=\! 9$. To stress-test heterogeneity, we
construct $K \!=\! 30$ synthetic clients by drawing label proportions from a
$\mathrm{Dir}(\beta\!\cdot\!\mathbf{1}_2)$ for $\beta \!\in\! \{0.1, 0.5, 1.0,
\infty\}$ ($\beta\!=\!\infty$ means uniform = pure IID). Each synthetic client
gets approximately $|\mathcal{D}|/K$ samples drawn from the global pool with
the assigned label distribution. We run \fedavg{}, \fedprox{}-$\mu\!=\!0.001$,
and \cvar{}-$\alpha\!=\!0.9$ across $10$ seeds for each $\beta$.

\textbf{Why-it-was-chosen.} This is the cleanest knob for heterogeneity that
exists. By varying only $\beta$ we can read off the heterogeneity-vs-method
slope and compare it to the within-method axes (Section~\ref{sec:mechanism-dirichlet}).

\subsection{Loss-landscape interpolation}
Following~[Li et al., 2018], we save three model weight vectors:
$w_{\text{init}}$ (random init), $w_{\text{best}}$ (lowest validation loss
checkpoint), $w_{\text{final}}$ (last round). For 1D, we sweep
$w_\alpha = (1-\alpha)\, w_{\text{init}} + \alpha\, w_{\text{final}}$ over
$\alpha \in [-0.5, 1.5]$ at $41$ points. For 2D, we sweep two random
filter-normalized directions on a $25\!\times\!25$ grid. Validation set is
$5{,}000$ stratified rows
(\texttt{outputs/full\_mimic\_iv\_proposal\_alignment/landscape/loss\_landscape\_config.json}).
Both backbones (logistic and MLP) get the full landscape.

\textbf{Why-it-was-chosen.} The landscape tells us whether the optimization
landscape is smooth (then DE/grid search differences are explained by mesh
density) or rough (then DE wins because grid hits non-minimum mesh points).
Section~\ref{sec:mechanism-landscape} reads this off the figures.
